% GENERATED FILE. Edit canonical lessons, then run npm run generate:books.
% canonical-source-hash: fnv1a64:da413ae9f8e2d40a
% canonical-lessons: MR-C329-talvar, MR-C329-dhal, MR-C329-sastra, MR-C329-toph, MR-C329-ghoddal

\chapter{Swords and Cannons}
\label{ch:mr-swords-and-cannons}
\hlchaptermodality{329}

\begin{chapteropening}
\textbf{By the end of this chapter:} \emph{I can name old weapons and the horsemen who carried them.}

Complete the last lesson of chapter 329: I can name old weapons and the horsemen who carried them.
\end{chapteropening}

\section[\texorpdfstring{\mr{तलवार} (talvār)}{तलवार (talvār)}]{\mr{तलवार} (talvār) — a sword}

\label{lesson:MR-C329-talvar}



\begin{quote}
Before the new one: say the Marathi for to fight, then the Marathi for a battlefield.
\end{quote}

\subsection*{You'll want to know: \mr{तलवार}}
\textbf{\mr{तलवार}} — \emph{talvār} — \textquotedblleft{}a sword\textquotedblright{}.

\begin{quote}
\mr{संग्रहालयात} \mr{एक} \mr{जुनी} \mr{तलवार} \mr{ठेवली} \mr{आहे}.
\end{quote}

\begin{quote}
\emph{saṅgrahālayāt ek junī talvār ṭhevlī āhe.}
\end{quote}

\begin{quote}
An old sword is kept in the museum.
\end{quote}

Feminine.

Plural \mr{तलवारी}.

\begin{cousinweb}
\textbf{\mr{तलवार}} (\emph{talvār}) is a tadbhava: CDIAL lists Marathi \emph{tarvār}, 'a one-edged sword', under Sanskrit \emph{taravāri} (CDIAL 5706); standard Marathi writes it with l for r.
\end{cousinweb}

\begin{grammarlens}[title={What you've built}]
One more word for telling what happened.
\end{grammarlens}

\subsection*{Your turn}
\begin{itemize}
\raggedright
  \item \emph{Say it:} \emph{talvār}
  \item \emph{Say it:} \emph{talvār}, once more
  \item \emph{Say it:} \emph{saṅgrahālayāt ek junī talvār ṭhevlī āhe.} — An old sword is kept in the museum.
  \item \emph{Recall:} say the Marathi for to fight, then the Marathi for a battlefield, then say \emph{talvār} again
\end{itemize}

\begin{tcolorbox}[breakable,colback=teal!4,colframe=teal!35!black,title={Before you move on}]
What does \mr{तलवार} mean? (\textquotedblleft{}A sword\textquotedblright{}.) Say it once more.
\end{tcolorbox}
\section[\texorpdfstring{\mr{ढाल} (ḍhāl)}{ढाल (ḍhāl)}]{\mr{ढाल} (ḍhāl) — a round shield; a shield}

\label{lesson:MR-C329-dhal}



\begin{quote}
Before the new one: say the Marathi for a battlefield, then the Marathi for a sword.
\end{quote}

\subsection*{You'll want to know: \mr{ढाल}}
\textbf{\mr{ढाल}} — \emph{ḍhāl} — \textquotedblleft{}a round shield; a shield\textquotedblright{}.

\begin{quote}
\mr{सैनिकाने} \mr{ढालीने} \mr{वार} \mr{अडवला}.
\end{quote}

\begin{quote}
\emph{sainikāne ḍhālīne vār aḍavlā.}
\end{quote}

\begin{quote}
The soldier blocked the blow with his shield.
\end{quote}

Feminine.

\mr{ढालीने}, with the shield.

\begin{cousinweb}
\textbf{\mr{ढाल}} (\emph{ḍhāl}) is a tadbhava: CDIAL lists Marathi \emph{ḍhāl}, 'a shield', under Sanskrit \emph{ḍhāla}, 'a shield' (CDIAL 5583).
\end{cousinweb}

\begin{grammarlens}[title={What you've built}]
One more word for telling what happened.
\end{grammarlens}

\subsection*{Your turn}
\begin{itemize}
\raggedright
  \item \emph{Say it:} \emph{ḍhāl}
  \item \emph{Say it:} \emph{ḍhāl}, once more
  \item \emph{Say it:} \emph{sainikāne ḍhālīne vār aḍavlā.} — The soldier blocked the blow with his shield.
  \item \emph{Recall:} say the Marathi for a battlefield, then the Marathi for a sword, then say \emph{ḍhāl} again
\end{itemize}

\begin{tcolorbox}[breakable,colback=teal!4,colframe=teal!35!black,title={Before you move on}]
What does \mr{ढाल} mean? (\textquotedblleft{}A round shield; a shield\textquotedblright{}.) Say it once more.
\end{tcolorbox}
\section[\texorpdfstring{\mr{शस्त्र} (śastra)}{शस्त्र (śastra)}]{\mr{शस्त्र} (śastra) — a weapon}

\label{lesson:MR-C329-sastra}



\begin{quote}
Before the new one: say the Marathi for a sword, then the Marathi for a round shield.
\end{quote}

\subsection*{You'll want to know: \mr{शस्त्र}}
\textbf{\mr{शस्त्र}} — \emph{śastra} — \textquotedblleft{}a weapon\textquotedblright{}.

\begin{quote}
\mr{शत्रूने} \mr{सगळी} \mr{शस्त्रं} \mr{खाली} \mr{ठेवली}.
\end{quote}

\begin{quote}
\emph{śatrūne sagḷī śastraṁ khālī ṭhevlī.}
\end{quote}

\begin{quote}
The enemy laid down all their weapons.
\end{quote}

Neuter.

Plural \mr{शस्त्रं}; \mr{शस्त्र} \mr{खाली} \mr{ठेवणे}, to lay down arms.

\begin{cousinweb}
\textbf{\mr{शस्त्र}} (\emph{śastra}) is a tatsama, taken from Sanskrit \emph{śastra}, 'a knife, a sword, a weapon' (Monier-Williams).
\end{cousinweb}

\begin{grammarlens}[title={What you've built}]
One more word for telling what happened.
\end{grammarlens}

\subsection*{Your turn}
\begin{itemize}
\raggedright
  \item \emph{Say it:} \emph{śastra}
  \item \emph{Say it:} \emph{śastra}, once more
  \item \emph{Say it:} \emph{śatrūne sagḷī śastraṁ khālī ṭhevlī.} — The enemy laid down all their weapons.
  \item \emph{Recall:} say the Marathi for a sword, then the Marathi for a round shield, then say \emph{śastra} again
\end{itemize}

\begin{tcolorbox}[breakable,colback=teal!4,colframe=teal!35!black,title={Before you move on}]
What does \mr{शस्त्र} mean? (\textquotedblleft{}A weapon\textquotedblright{}.) Say it once more.
\end{tcolorbox}
\section[\texorpdfstring{\mr{तोफ} (toph)}{तोफ (toph)}]{\mr{तोफ} (toph) — a cannon}

\label{lesson:MR-C329-toph}



\begin{quote}
Before the new one: say the Marathi for a round shield, then the Marathi for a weapon.
\end{quote}

\subsection*{You'll want to know: \mr{तोफ}}
\textbf{\mr{तोफ}} — \emph{toph} — \textquotedblleft{}a cannon\textquotedblright{}.

\begin{quote}
\mr{किल्ल्याच्या} \mr{दारात} \mr{एक} \mr{मोठी} \mr{तोफ} \mr{आहे}.
\end{quote}

\begin{quote}
\emph{killyācyā dārāt ek moṭhī toph āhe.}
\end{quote}

\begin{quote}
There is a big cannon at the fort's gate.
\end{quote}

Feminine.

Plural \mr{तोफा}.

\begin{cousinweb}
\textbf{\mr{तोफ}} (\emph{toph}) is a loan from Persian \emph{top}, 'a cannon', usually traced to Turkish.
\end{cousinweb}

\begin{grammarlens}[title={What you've built}]
One more word for telling what happened.
\end{grammarlens}

\subsection*{Your turn}
\begin{itemize}
\raggedright
  \item \emph{Say it:} \emph{toph}
  \item \emph{Say it:} \emph{toph}, once more
  \item \emph{Say it:} \emph{killyācyā dārāt ek moṭhī toph āhe.} — There is a big cannon at the fort's gate.
  \item \emph{Recall:} say the Marathi for a round shield, then the Marathi for a weapon, then say \emph{toph} again
\end{itemize}

\begin{tcolorbox}[breakable,colback=teal!4,colframe=teal!35!black,title={Before you move on}]
What does \mr{तोफ} mean? (\textquotedblleft{}A cannon\textquotedblright{}.) Say it once more.
\end{tcolorbox}
\section[\texorpdfstring{\mr{घोडदळ} (ghoḍdaḷ)}{घोडदळ (ghoḍdaḷ)}]{\mr{घोडदळ} (ghoḍdaḷ) — cavalry}

\label{lesson:MR-C329-ghoddal}



\begin{quote}
Before the new one: say the Marathi for a weapon, then the Marathi for a cannon.
\end{quote}

\subsection*{You'll want to know: \mr{घोडदळ}}
\textbf{\mr{घोडदळ}} — \emph{ghoḍdaḷ} — \textquotedblleft{}cavalry\textquotedblright{}.

\begin{quote}
\mr{घोडदळ} \mr{वेगाने} \mr{पुढे} \mr{गेलं}.
\end{quote}

\begin{quote}
\emph{ghoḍdaḷ vegāne puḍhe gelaṁ.}
\end{quote}

\begin{quote}
The cavalry went forward at speed.
\end{quote}

Neuter.

\mr{घोडा}, a horse, + \mr{दळ}, a troop.

\begin{cousinweb}
\textbf{\mr{घोडदळ}} (\emph{ghoḍdaḷ}) is built from \emph{ghoḍā}, 'a horse', and \emph{daḷ}, 'a troop', from Sanskrit \emph{dala}, 'a part, a troop' (Monier-Williams).
\end{cousinweb}

\begin{grammarlens}[title={What you've built}]
One more word for telling what happened.
\end{grammarlens}

\subsection*{Your turn}
\begin{itemize}
\raggedright
  \item \emph{Say it:} \emph{ghoḍdaḷ}
  \item \emph{Say it:} \emph{ghoḍdaḷ}, once more
  \item \emph{Say it:} \emph{ghoḍdaḷ vegāne puḍhe gelaṁ.} — The cavalry went forward at speed.
  \item \emph{Recall:} say the Marathi for a weapon, then the Marathi for a cannon, then say \emph{ghoḍdaḷ} again
\end{itemize}

\begin{tcolorbox}[breakable,colback=teal!4,colframe=teal!35!black,title={Before you move on}]
What does \mr{घोडदळ} mean? (\textquotedblleft{}Cavalry\textquotedblright{}.) Say it once more.
\end{tcolorbox}
